\documentclass[mlabstract,onecolumn]{jmlr}

\usepackage{booktabs}
\usepackage{multirow}
\usepackage{amsmath}
\usepackage{capt-of}

\newcommand{\auc}{\mathrm{AUC}}

%%%% DON'T CHANGE %%%%%%%%%
\jmlrvolume{}
\firstpageno{1}
\editors{}
\jmlryear{2026}
\jmlrworkshop{Symmetry and Geometry in Neural Representations}
%%%%%%%%%%%%%%%%%%%%%%%%%%%

\title[Cortical Decoders Age Slower]{Cortical Decoders Trained on Other
Individuals Age Slower Than Your Own}

\author{\Name{Anonymous Author(s)}}

\begin{document}
\maketitle

\begin{abstract}
With adapting brain--computer interface (BCI) technology, neural decoding algorithms are one of the
most essential parts of being able to understand people's brains. Decoders must work for months, and
recalibrating them per subject has high subject costs. A decoder fit to a singular animal's cortex
drifts, losing accuracy over days, and usually the remedy is recalibrating on fresh data from that
same subject. We report a different kind of robustness: a
decoder built from one other individual, or pooled from a set of other individuals, ages more slowly
than the subject's own. On the open BraiDyn-BC cohort (25 mice, 44 Allen-atlas cortical parcels, a
$\sim$15-day operant protocol) we decode four task events and measure, per mouse, the $\auc$ lost
over two weeks by its own early-week decoder against one pooled over the other 24 mice and one built
from a single other mouse. The personal decoder ages more by $+0.017$ $\auc$
($p = 5 \times 10^{-5}$, 73 of 100 mouse--event pairs). This also is not a data-volume effect: at a matched training-event count a
decoder from a single other mouse still ages less. Hence we conclude this is rather an effect of not
overfitting an individual's drift-prone idiosyncrasies. It also survives a 1-D CNN and a GRU. The
effect holds past this cohort and recording modality. On an independent two-photon dataset from a
different laboratory -- Allen Visual Behavior, 23 mice, single-plane VISp populations decoding a
stimulus change -- the same session-held-out protocol gives $+0.023$ $\auc$
($p = 0.003$, 74\% of mice), and the count-matched control reproduces there as well
(self $-$ one other mouse $= +0.011$ $\auc$). The
early-week accuracy cost of decoding across individuals is repaid as late-week stability, which
suggests population priors as an alternative to per-subject recalibration. A second offering is the
session-held-out protocol under which all of this is measured: because only the within-subject arm
can share a recording session with its test data, a trial-level split inflates that arm alone, and
holding out whole sessions is what allows this form of robustness to be reported reliably. Code and
a streamed-feature pipeline are released.
\end{abstract}

\section{Introduction}
Neural decoding is one of the most essential ways to understand neural interfaces and create BCI
technology, and a decoder that is deployed is a decoder that has to keep working: over months, on
the same individual, without fresh labelled data arriving to prop it up. Representational drift
stands directly against that. The reorganization of an animal's neural code over days with fixed
behavior is now documented across sensory, motor, and association cortex and in the hippocampus
\citep{rule2020,driscoll2017,schoonover2021}. This has a stark impact on neural decoding: a decoder
trained on an animal today degrades on that same animal weeks later. The dominant response is
recalibration, periodically refitting the decoder on fresh within-subject data
\citep{degenhart2020,jarosiewicz2015}, which requires ongoing labelled data from the individual and
treats each subject in isolation.

The code's cross-individual structure offers an alternate route to temporal robustness, and the
field already relies on it. Cortex-wide behavioral representations are considerably shared among
individuals \citep{musall2019,macdowell2020,macdowell2024}; latent dynamics are preserved across
animals well enough that a decoder fit to one transfers to another \citep{safaie2023}; and pooled
decoders already read out behavior from unseen individuals \citep{wicat2026,neds2025,peterson2021}.
Methods such as Meta-AlignNN \citep{metaalign2025} are built on that premise. While people currently
know that conservation across individuals is established, and that pooling is known to work at a
fixed time, every one of these results is measured at a fixed time. We offer pooling as a solution
to a new problem: the fact that it resists drift over time.

Our core discovery is that a decoder that never saw the subject cannot fit to that subject's
drift-prone idiosyncrasies, so it should age more slowly, and the way we measure that difference is
as an asymmetry between the two arms. This prediction, to our knowledge, has not been examined.

We test it on BraiDyn-BC \citep{kondo2025}, a widefield calcium-imaging resource whose $\sim$15-day
operant protocol and shared 44-region Allen parcellation make within-subject and cross-subject
decoders directly comparable across time. Our contributions:
\begin{itemize}\itemsep2pt
\item \textbf{Pooling resists drift.} Across a two-week period the population-pooled decoder loses
less accuracy than the individual's own, with a paired asymmetry of $+0.017$ $\auc$
($p = 5 \times 10^{-5}$, 73 of 100 pairs). The size of this effect varies
among the four events, but those are chained within a trial and have highly differing counts, so we
treat the pooled estimate as the result and leave the per-event ranking uninterpreted.
\item \textbf{It is not a data-volume artifact.} At a matched training-event count a decoder trained
on one other mouse ages less than one trained on the subject's own, by $+0.010$ $\auc$ on average.
The gain therefore comes from training on individuals other than the subject, not from the larger
training set pooling usually brings, and training on many other mice rather than one lowers aging
further on all four events.
\item \textbf{It replicates on an independent dataset.} On the Allen Visual Behavior two-photon
cohort the same protocol gives $+0.023$ $\auc$ ($p = 0.003$, 74\% of 23 mice), with the
count-matched ordering also reproduced.
\item \textbf{Session-level validation is load-bearing.} As trials within a recording session are
dependent, a trial-level split inflates the personal decoder but not the pooled one, one-sidedly
inflating the asymmetry. We report that inflation for each event; on lever pull it accounts for the
entire apparent effect. For longitudinal decoding comparisons we recommend session-held-out
validation as a default.
\end{itemize}



\section{Data and methods}
BraiDyn-BC \citep{kondo2025} (DANDI:001425) provides dorsal-cortex $\Delta F/F$ parcellated into 44
Allen Common Coordinate Framework (CCF) regions for 25 mice, each completing $\sim$15 daily operant
sessions over $\sim$3 weeks. For each onset of four events (lever pull, lick, reward, tone) we build
a 44-dimensional evoked feature and decode it against matched negatives with $\ell_2$-regularized
logistic regression, scored by area under the receiver-operating characteristic curve ($\auc$).
Sessions are grouped into an early block (days 1--5) and a late block (days 11--15). The personal
decoder leaves out one whole early \emph{session} at a time; the pooled decoder trains on the other
mice's early blocks. Aging is the accuracy lost as testing moves from early to late with training
fixed, and the per-mouse asymmetry is personal aging minus pooled aging, positive when the personal
decoder ages more. Significance is a paired sign-flip permutation test across mice
(20{,}000 permutations), with cluster-bootstrap confidence intervals. Holding out whole sessions
rather than trials is essential rather than incidental: only the within-subject arm can share a
session between train and test, so a trial-level split inflates that arm alone.
The external validation cohort is Allen Visual Behavior two-photon \citep{allenvb2023}, for which
the absence of any neuron correspondence across animals forces a permutation-invariant 27-dimensional
summary of the population response. Appendix~\ref{app:methods} gives both in full.

\section{Results}
For each mouse we keep the test data fixed (first the mouse's held-out early sessions, then its late
block) and change only the source of the training data: either the mouse itself or the other 24 mice.
Across all 100 mouse--event pairs the asymmetry is $+0.017$ $\auc$ ($p = 5 \times 10^{-5}$), with the
personal decoder aging more in 73\% of pairs; the per-event breakdown is in Table~\ref{tab:asym}.

The effect differs among the four events: it is largest for reward ($+0.029$, $p = 0.0005$) and lick
($+0.020$, $p = 0.002$), marginal for tone ($+0.015$, $p = 0.08$), and absent for lever pull
($+0.003$, $p = 0.37$). While we report this variation, we do not interpret it, as the four events
all fall within roughly the same two seconds of a trial. Accordingly, the pooled estimate over all
100 pairs is the claim we defend, and the per-event column should be read as data rather than as
four separate findings.



A pooled decoder contains $\sim$24$\times$ more training events than a personal decoder, so its
stability could reflect data volume rather than pooling. We test this with a count-matched control
For each held-out early session $s$ of each mouse we set the
training-set size to $n$, the number of events in that mouse's remaining early sessions, and compare
three decoders: \emph{self}, trained on those $n$ events; \emph{one}, trained on $n$ events from a
single other mouse; and \emph{many}, trained on $n$ events drawn across the other 24 mice. All three
are scored on the same held-out session $s$ and the same late block, so they differ only in whose
data they were fit on.

For the events where there was an asymmetry, the count-matched ordering reproduces it. The self
decoder ages more than a decoder trained on one other mouse by $+0.010$ $\auc$ on average, and by
$+0.007$ for lick, $+0.017$ for reward and $+0.006$ for tone individually. As both decoders see the
same number of training events, this is not a data-volume artifact.

To further test whether the asymmetry is a property of that one cohort, modality and task, we repeat
the analysis on the Allen Visual Behavior two-photon dataset \citep{allenvb2023}: a different
laboratory, cellular two-photon rather than mesoscale widefield imaging, a visual change-detection
task rather than a cued lever pull, and a 27-dimensional permutation-invariant population summary
rather than a 44-parcel atlas feature. The asymmetry reproduces at a comparable magnitude
(Table~\ref{tab:allen}): across 23 mice the personal decoder ages by $+0.024$ $\auc$ and the pooled
decoder by $+0.001$, an asymmetry of $+0.023$ $[0.010, 0.037]$, $p = 0.003$, with the mouse's own
decoder ageing more in 74\% of animals. The count-matched control reproduces as well, at $+0.011$
$[0.001, 0.022]$ for self minus one other mouse, though on 23 animals it is marginal by the sign-flip
test ($p = 0.054$). The mechanism is not identical in the two cohorts -- on BraiDyn-BC the gap
narrows because the pooled decoder gains, here because the personal decoder loses -- and the shared
finding is that the personal decoder ages faster. The inclusion-rule sweep and the decomposition of the trial-split behaviour on this cohort are
reported in the full paper.

\begin{table}[t]\centering
\caption{The aging asymmetry on the original cohort and on the external validation dataset. Aging
is $\auc$ lost from early to late test with training held fixed; asymmetry is personal $-$ pooled
aging, positive when the subject's own decoder ages more. Both rows use the same session-held-out
protocol and the same count-matched control. $p$: paired sign-flip permutation.}
\label{tab:allen}
\setlength{\tabcolsep}{3pt}
{\small
\begin{tabular}{lccccc}
\toprule
& Personal & Pooled & Asymmetry (95\% CI) & $p$ & own ages more \\
\midrule
BraiDyn-BC, widefield, 25 mice & & & $+0.017$ $[0.010, 0.024]$ & $5\times10^{-5}$ & 73\% \\
Allen VB, two-photon, 23 mice & $+0.024$ & $+0.001$ & $+0.023$ $[0.010, 0.037]$ & 0.003 & 74\% \\
\bottomrule
\end{tabular}}

\vspace{3pt}
\begin{minipage}{0.95\linewidth}\small\raggedright
\emph{Count-matched control, aging at equal training-event count.} BraiDyn-BC: self $-$ one other
mouse $= +0.010$; many others age less than one on all four events. Allen VB: self $+0.024$, one
$+0.013$, many $+0.003$; self $-$ one $= +0.011$, many $-$ one $= -0.010$.
\end{minipage}
\end{table}

\section{Discussion}
A subject's own decoder overfits idiosyncratic features of its early sessions, some of which later
drift. A decoder that has never seen the subject must instead rely on components of the code that are
shared across individuals and more stable over time. The count-matched control supports this reading
directly: at equal data, other-individual data ages less than the subject's own. Unlike
within-subject recalibration, this requires no newly labelled data from the individual at test
time -- the practical point, since obtaining that data is the expensive part of deploying a decoder
over months.


The validation protocol is worth stating explicitly: a
trial-level split inflates the personal decoder's within-block accuracy by up to $+0.013$ $\auc$ and
leaves the pooled decoder untouched, and on lever pull that inflation accounts for the entire
apparent effect ($+0.014$, $p = 0.001$, becoming $+0.003$, $p = 0.37$). The practical consequences are collected as a checklist in Appendix~\ref{app:practice}.

\textbf{Limitations.} The four events are not four independent behaviors but four thresholded
channels of one cued lever-pull trial, so we cannot always say which behavior the effect belongs to.
We measure at the mesoscale, and the two-photon replication reaches single neurons only at one depth
in one visual area. The trial-level inflation above is a caution for comparisons of this shape
rather than a limitation of this work: every number we report is measured after holding out whole
sessions.



\begin{thebibliography}{99}
\bibitem[Allen Institute for Brain Science(2023)]{allenvb2023} Allen Institute for Brain Science.
Visual Behavior --- 2-photon. \emph{Allen Brain Observatory}, 2023.
\url{https://portal.brain-map.org/circuits-behavior/visual-behavior-2p}.
\bibitem[Degenhart et al.(2020)]{degenhart2020} A.~D.~Degenhart et al. Stabilization of a
brain--computer interface via the alignment of low-dimensional spaces of neural activity.
\emph{Nature Biomedical Engineering}, 4:672--685, 2020.
\bibitem[Driscoll et al.(2017)]{driscoll2017} L.~N.~Driscoll et al. Dynamic reorganization of
neuronal activity patterns in parietal cortex. \emph{Cell}, 170(5):986--999, 2017.
\bibitem[Farshchian et al.(2019)]{farshchian2019} A.~Farshchian et al. Adversarial domain
adaptation for stable brain--machine interfaces. \emph{ICLR}, 2019.
\bibitem[Jarosiewicz et al.(2015)]{jarosiewicz2015} B.~Jarosiewicz et al. Virtual typing by people
with tetraplegia using a self-calibrating intracortical brain--computer interface. \emph{Science
Translational Medicine}, 7(313):313ra179, 2015.
\bibitem[Kondo et al.(2025)]{kondo2025} A.~Kondo et al. BraiDyn-BC: a cortex-wide widefield calcium
imaging dataset during a cued lever-pull operant task. \emph{Scientific Data}, 12, 2025.
DANDI:001425.
\bibitem[MacDowell \& Buschman(2020)]{macdowell2020} C.~J.~MacDowell and T.~J.~Buschman.
Low-dimensional spatiotemporal dynamics underlie cortex-wide neural activity. \emph{Current
Biology}, 30(14):2665--2680, 2020.
\bibitem[MacDowell et al.(2024)]{macdowell2024} C.~J.~MacDowell et al. Shared cortex-wide
spatiotemporal motifs govern activity across individuals and behavioral states. \emph{Current
Biology}, 34, 2024.
\bibitem[Melbaum et al.(2022)]{melbaum2022} S.~Melbaum et al. Conserved structures of neural
activity in sensorimotor cortex of freely moving rats enable cross-subject decoding. \emph{Nature
Communications}, 13:7902, 2022.
\bibitem[Zou et al.(2025)]{metaalign2025} Y.~Zou, K.~Liu, C.~Zhang, Y.~Ling, F.~Wang, M.~Li,
Y.~Chen, M.~Li, S.~Guan, Z.~He, and C.~T.~Li. Meta-AlignNN: a meta-learning framework for stable
brain--computer interface performance across subjects, time, and tasks. \emph{bioRxiv}, 2025.
doi:10.1101/2025.04.20.649482.
\bibitem[Musall et al.(2019)]{musall2019} S.~Musall et al. Single-trial neural dynamics are
dominated by richly varied movements. \emph{Nature Neuroscience}, 22(10):1677--1686, 2019.
\bibitem[Zhang et al.(2025)]{neds2025} Y.~Zhang, Y.~Wang, M.~Azabou, A.~Andre, Z.~Wang, H.~Lyu,
The International Brain Laboratory, E.~Dyer, L.~Paninski, and C.~Hurwitz. Neural encoding and
decoding at scale. arXiv:2504.08201, 2025.
\bibitem[Peterson et al.(2021)]{peterson2021} S.~M.~Peterson, Z.~Steine-Hanson, N.~Davis,
R.~P.~N.~Rao, and B.~W.~Brunton. Generalized neural decoders for transfer learning across
participants and recording modalities. \emph{Journal of Neural Engineering}, 18:026014, 2021.
\bibitem[Karpowicz et al.(2025)]{karpowicz2025} B.~M.~Karpowicz, Y.~H.~Ali, L.~N.~Wimalasena,
et al. Stabilizing brain--computer interfaces through alignment of latent dynamics (NoMAD).
\emph{Nature Communications}, 16:4662, 2025.
\bibitem[Wilson et al.(2025)]{wilson2025} G.~H.~Wilson et al. Long-term unsupervised recalibration
of cursor-based intracortical brain--computer interfaces using a hidden Markov model. \emph{Nature
Biomedical Engineering}, 2025.
\bibitem[Rule et al.(2019)]{rule2019} M.~E.~Rule, T.~O'Leary, and C.~D.~Harvey. Causes and
consequences of representational drift. \emph{Current Opinion in Neurobiology}, 58:141--147, 2019.
\bibitem[Safaie et al.(2023)]{safaie2023} M.~Safaie, J.~C.~Chang, J.~Park, L.~E.~Miller,
J.~T.~Dudman, M.~G.~Perich, and J.~A.~Gallego. Preserved neural dynamics across animals performing
similar behaviour. \emph{Nature}, 623:765--771, 2023.
\bibitem[Rule et al.(2020)]{rule2020} M.~E.~Rule et al. Stable task information from an unstable
neural population. \emph{eLife}, 9:e51121, 2020.
\bibitem[Schneider et al.(2023)]{schneider2023} S.~Schneider, J.~H.~Lee, and M.~W.~Mathis.
Learnable latent embeddings for joint behavioural and neural analysis (CEBRA). \emph{Nature},
617:360--368, 2023.
\bibitem[Schoonover et al.(2021)]{schoonover2021} C.~E.~Schoonover et al. Representational drift in
primary olfactory cortex. \emph{Nature}, 594:541--546, 2021.
\bibitem[Hosseini et al.(2026)]{wicat2026} M.~Hosseini, E.~Erturk, S.~Hashemi, and M.~M.~Shanechi.
Cross-subject modeling for widefield calcium imaging via atlas-aligned spatiotemporal tokenization.
In \emph{Proceedings of the 43rd International Conference on Machine Learning (ICML)}, 2026.
arXiv:2607.09754.
\end{thebibliography}

\appendix

\section{Best-practice checklist}
\label{app:practice}
\begin{center}
\fbox{\begin{minipage}{0.95\linewidth}
\small
\textbf{Best practice for longitudinal within- versus cross-subject comparisons}
\begin{enumerate}\itemsep0pt\parskip0pt\topsep2pt
\item \textbf{Hold out whole recording sessions, not trials.} Only the within-subject arm can share
a session between train and test. Here a trial-level split inflated it by up to $+0.013$ $\auc$ and
on lever pull accounted for the entire apparent effect ($+0.014$, $p=0.001 \to +0.003$, $p=0.37$).
\item \textbf{Count-match the training sets, and never score an arm on the block it was fit on.}
A pooled decoder holds $\sim$24$\times$ more events; scoring \emph{self} on its own training block
is not a cross-validated quantity and reverses the ordering the control exists to establish.
\item \textbf{Average per-session scores rather than pooling one score across sessions,} and report
each arm's aging separately. On the Allen cohort a single pooled $\auc$ cost $-0.030$
($p<10^{-4}$) and dominated the leakage term it was confounded with.
\end{enumerate}
\end{minipage}}
\captionof{table}{}\label{box:practice}
\end{center}

\section{Per-event results}
\label{app:perevent}
\begin{table}[t]\centering
\caption{Personal vs.\ pooled aging across days ($\auc$ lost from early to late test, holding
training fixed), per event, $n = 25$ mice, with whole early sessions held out. Asymmetry $=$
personal $-$ pooled aging; a positive value means the mouse's own decoder ages more. $p$: paired
sign-flip permutation. The rightmost column gives the same asymmetry under a trial-level split, the
protocol this analysis replaces.}
\label{tab:asym}
\setlength{\tabcolsep}{4pt}
\begin{tabular}{lccccccc}
\toprule
Event & Personal & Pooled & Asymmetry (95\% CI) & $p$ & own ages more & trial-split \\
\midrule
Lever-pull & $-0.022$ & $-0.025$ & $+0.003$ $[-0.004, 0.011]$ & 0.37 & 64\% & $+0.014$ \\
Lick       & $+0.012$ & $-0.009$ & $+0.020$ $[0.008, 0.032]$ & 0.0023 & 80\% & $+0.028$ \\
Reward     & $+0.017$ & $-0.012$ & $+0.029$ $[0.014, 0.046]$ & 0.0005 & 80\% & $+0.032$ \\
Tone       & $+0.016$ & $+0.001$ & $+0.015$ $[0.001, 0.032]$ & 0.077 & 68\% & $+0.022$ \\
\midrule
\textbf{Pooled} & & & $\mathbf{+0.017}$ $[0.010, 0.024]$ & $\mathbf{5\times10^{-5}}$ & \textbf{73\%} & $+0.024$ \\
\bottomrule
\end{tabular}
\end{table}

\section{Data and methods in full}
\label{app:methods}
\subsection{Datasets and preprocessing}
BraiDyn-BC \citep{kondo2025} (DANDI:001425, CC-BY) provides, for each mouse and
session, hemodynamically corrected dorsal-cortex $\Delta F/F$ traces parcellated into 44 Allen Common
Coordinate Framework regions at $\sim$30\,Hz, together with time-aligned behavioral channels. The
dataset includes 25 mice, each of which completed $\sim$15 daily operant sessions over $\sim$3 weeks.
We stream the 44-region traces and behavioral channels without downloading the raw movies.

The two cohorts used in this paper were built for different purposes, and the ways they differ are
what make the comparison informative rather than redundant. BraiDyn-BC is densely longitudinal, with
roughly fifteen sessions per mouse over three weeks, and it supplies a shared coordinate system: the
44 Allen parcels mean that region $k$ in one mouse is the same region in every other, so a pooled
decoder can be fit directly in a common feature space. The Allen Visual Behavior cohort has neither
property. Its longitudinal subset is sparser, and because it is cellular two-photon imaging there is
no correspondence between animals at all -- one mouse's cell 17 is unrelated to another's -- so a
common feature space has to be constructed rather than assumed (Section~\ref{sec:allencohort}). The
recording modality, the task, and the decoded contrast all differ as well. A result that survives
both is therefore not a property of one atlas, one instrument, or one behavior.

\subsection{Features and decoders}
We detect the onset of four events (lever pull, lick, reward, and
tone) as rising edges, with a 1\,s refractory period. For each onset we construct a 44-dimensional
evoked feature by subtracting the mean pre-onset baseline ($-0.5$--$0$\,s) from the mean post-onset
$\Delta F/F$ response ($0$--$1$\,s) in each parcel. Negative samples are drawn from matched-count
windows more than 2\,s from any onset \emph{of the same event}; they are not screened against the
other three events. Every
classification is a binary decode of one event against these matched negatives, scored by area under
the ROC curve ($\auc$). The four decoders are therefore four independent binary detectors, not a
four-way classifier: nothing in this paper asks a model to distinguish one event from another.

The primary decoder is a standardized $\ell_2$-regularized logistic regression on
the 44-dimensional evoked feature. This is a linear readout of the parcel-space code: it recovers the
component of behavior that is linearly present in the population average, and it is the natural
comparison to prior cross-subject decoders, which are overwhelmingly linear
\citep{melbaum2022,peterson2021}. Every condition below uses this same readout, so the only variable
is which data the decoder was trained on.\footnote{The result is not a property of the linear
readout. Repeating the analysis with a 1-D convolutional network over time and a GRU over the frame
sequence (full evoked window, $-0.5$ to $+1.0$\,s, $45 \times 44$), under the same session-held-out
protocol, the asymmetry survives for both events that carry it: lick gives $+0.020$, $+0.012$ and
$+0.010$ and reward $+0.029$, $+0.015$ and $+0.012$ under the linear, CNN and GRU readouts, all
significant. Under a trial-level split the nonlinear estimates appear \emph{larger} instead, which is
the trial-level inflation: a higher-capacity model exploits it more effectively
than a linear one.}

\subsection{Temporal blocks, aging, and the session-held-out protocol}
\label{sec:blocks}
We group each mouse's sessions into an early block (days 1--5)
and a late block (days 11--15) and report $\auc$ under four conditions. The personal decoder is
estimated by leaving out one whole early \emph{session} at a time: for each early session $s$ we
train on the mouse's remaining early sessions and score the resulting decoder both on $s$
(within-mouse same-block, WS) and on the late block (within-mouse across-block, WX), then average
over $s$. Because WS and WX come from the same fitted decoders, they differ only in which data they
are tested on. The pooled decoder is trained once on the other mice's early blocks and scored on the
same held-out sessions $s$ (LS) and on the late block (LX).

Holding out sessions rather than trials is essential and not a detail. Trials within one recording
share illumination and hemodynamic state, slow arousal drift, and temporal autocorrelation, so a
trial-level split lets the personal decoder train and test inside the same session. The pooled
decoder is trained on other animals entirely and cannot benefit this way, so the inflation is
one-sided and lands directly on the quantity of interest; the main text reports what it
costs per event.

Aging is the accuracy lost when testing moves from the held-out mouse's early data to its late block
while training stays fixed: personal aging is $\mathrm{WS}-\mathrm{WX}$ and pooled aging is
$\mathrm{LS}-\mathrm{LX}$. The per-mouse aging asymmetry is
$(\mathrm{WS}-\mathrm{WX})-(\mathrm{LS}-\mathrm{LX})$; a positive value means the personal decoder
ages more. We assess significance with a paired sign-flip permutation test across mice
(20{,}000 permutations) and compute 95\% confidence intervals with a cluster bootstrap over mice.

\subsection{The external validation cohort}
\label{sec:allencohort}
The external validation repeats the analysis on the Allen
Visual Behavior two-photon dataset \citep{allenvb2023}. We select containers spanning at least 14
days whose experiments are all in VISp within a single 50\,\textmu m depth bin, which is the
cross-subject correspondence available here; the largest such band, 150--199\,\textmu m, holds 25
mice and 184 experiments. We require at least 15 segmented cells and 40 usable trials per
experiment, leaving 166 experiments and 23 mice with at least four sessions. Sessions are streamed
from S3 rather than downloaded.

As different animals have no neuron correspondence and the longitudinal subset is single-plane,
averaging over the field of view would give a one-dimensional feature. The representation that is
comparable across animals is
therefore a permutation-invariant summary of the population response: per neuron we take the mean
$\Delta F/F$ over $0$--$1$\,s minus the mean over $-0.5$--$0$\,s, then describe the distribution of
those evoked responses by 25 quantiles plus its mean and standard deviation, giving a fixed
27-dimensional feature regardless of how many cells were segmented. The decoded contrast is
\texttt{is\_change}, a real image change against a sham catch trial, which is a stimulus contrast
rather than a reward- or lick-coupled one. Each mouse's sessions are split by date into early and
late halves, and the personal decoder, the pooled decoder, aging, the asymmetry and the
count-matched control are all defined exactly as above.


\end{document}
